%*********************************************************************
%	CONCLUSIONS
%*********************************************************************
\chapter{Conclusions}
\label{ch:conclusions}

The goal of this project was to investigate and improve upon the current system in place used to take depth scans of intertidal sediment, for the purpose of surface rugosity metrics. The main investigation behind these metrics is trying to understand the roll that human activity effects the biodensity of the benthic ecosystem within intertidal regions at the Langebaan Lagoon. It has been identified in relating literature, that one metric that can be used to estimate this biological activity is surface rugosity. Biological activity in these benthic ecosystems are shown to be strong bioturbators, and have significant impact on many aspects of the sediment. One such aspect is the surface, where activity, such as burrowing, disturb the natural sediment, and inflate the surface roughness. Therefore more activity should have a correlation with higher surface roughness. \todo{maybe site some references here}.

The project as it was before this investigation was based off that of \citet{azharRGBDFrameworkCapturing2022}, whereby a reletivly inexpensive RGB-D camera was used to get surface reconstrucitons of the sediment. There are many benefits to this approach, mainly being the non contact based measurement of the surface, which prevents disturbing the ecosystem that is being investigated, during measurement. Secondly, the 3D reconstruction provides a more accurate measure of surface roughness, since the reconstruction provides a three dimensional surface, rather than a single dimension estimate, commonly used in previous investigations. Where the two projects differentiate is the use of camera technology. There are several potential ways of obtaining three dimensional reconstructions of surface, including LiDAR, Stereo vision and SFM. These methods have significant downsides to them, Mainly high cost, for example using LiDAR technology, or high computational requirements for actual reconstruction, for example using traditional stereo vision, or SFM. This project uses an Intel RealSense D456 active stereo vision camera, that provides an easy way to obtain surface scans, thanks to its onboard processing, which provide the user with premade depth maps. This camera was used instead of the ASUS one used by \citet{azharRGBDFrameworkCapturing2022}, since it features water resistant ratings, and uses active stereo vision, rather than structre light, allowing it to operate under direct sun exposure.

One aim of this project was to investigate the performance of the D456, and to develop further processing pipelines to improve the reconstructions. The first investigation conducted was to characterise the performance of the Intel RealSense D456 under a variety of conditions that are generally present in conditions similar to those present in intertidal regions. This includes harsh lighting conditions such as direct sunlight, and variability of operating conditions such as height and tilt. This was done by evaluating its ability to reconstruct both planar and 3D surfaces at a variety of heights, angles lighting conditions, camera paramters, depth filtering, and fusion configuration. By inspecting the performance of the camera under these different conditions, the optimal camera configration can be found for this use case.

The second aim of this project was to investigate more post-processing steps that could be taken to further improve the accuracy of these surfaces scans. the main investigation here was developing a fusion pipeline that would take multi-view scans and fuse them together into a single higher quality reconstruction, that could in turn be used for analysis.


Based of the results shown in \autoref{ch:experiments} and \autoref{ch:Validation}, it can be seen that the main factors effecting the performance of the RealSense D456 camera are scan distance, and resolution. At farther distance scans, the reconstrucion both in terms of planar surfaces and 3D surfaces degrades with an approximate linear trend at this scale, showing that the closer the camera is to the surface being scaned, the higher the accuracy is. Additionally, at higher image resolutions, the camera is able to procduce higher quality depth scans. \todo{Word this nicer, and use some metrics} Other results show minimal or trade off results. The camera angle has a negative effect when representing planar surfaces, however produces better results on scans  for the 3D surface, this can be attributed to the fact that 3D surfaces feature some relief that can be hidden by the structure of the surface, and at different angles, this relief may be better represented. The application of filters, and preset changes, do have some impact on the surface, but three is a minimal change, and for the case with filters, these can be applied to the scans in post, and thus show that on site capturing would benefit from no filters being applied, to shorten the scan period. Depth units showed less improvement than expected, however consideration does need to be taken, as the reference objects used for this characterisation have large, quickly changing z heights, and so a more smooth, shallow height change might benefit from lower depth units, as the z quantisation would change.

The results from fusion, show strong improvements over single scans. It is generally seen that the more viewpoints taken, the higherr the overal result. But one factor to consider with these, is that an individual scancan take up to \todo{Find how long the scans take}, at full sweeps. This is impractical in situ, where up to 40 different scans need to be taken, in a small window low tide window period. From the results oit can be seen however, that taking fewer scans with larger increments produces similarly better results, and so using a combination of say 0$\pm$20\textdegree may provide results considered satisfactory.

Validation was performed, by inspecting the roughness metric output of these scans compared to both the known reference objects and sediment taken from the actual Langebaan Lagoon. These results do provide a more qualitative analysis, but provide useful insight into interpreting other results. The rugosity metrics on the reference objects, whilst initially seem to contradict the results of characterisation, do in fact show that fusion does reduces depth noise, and so better reflects the true conditions of the surface, which can be seen in the sand tests. Where the flat surfaces show that the fused results better represent the true state of the surface, compared to the single scans. \todo{Reference the bumpy single scans, and images, and rugosity results}

Some additional tools were made in aims to help the functionality of the processing pipeline. An automatic cropping tool was made, which crops each individual scan to the region bounded by four april tags present on the legs of the rig. The cropping tool keeps track of the order of the tags, and extraoplates their position in the case where one may not be detected. This allows the program to remain robust in the case of both misidentified tags, and tags not in view. An aditional function is that the code provides a reference point, after cropping the tags out. This reference point can be used by the fusion pipeline for coarse alignemnt, and allows for a higher accuracy alignment compared to other methods.

Overall the characterisation found useful results, which will help further projects when aiming to use the D456 camera for reconstruction. And the processing pipeline shows good results regarding the reconstruction accuracy of the system especially with multi-view scans.
%---------------------------------------------------------------------

%*********************************************************************